\section{The mass determined by the functional}\label{sec:normalization-mass-quotient}

For closed backgrounds, the functional has a precise loss of information: on the model domain below, equality of functionals is exactly equality of mass. This also gives examples in which supporting probabilities stay fixed while the full potentials have different limits.

\Needspace{9\baselineskip}
In this section $X$ is compact K\"ahler. Allow $\theta$ to be any smooth closed big form, and let $\varphi=P_\theta[\varphi]$ have small unbounded locus in the sense of \cref{sec:psh-singularity-types} and positive mass
\[
 m_\varphi=\int_X\mathrm{MA}_\theta(\varphi)>0.
\]
Write $B_{\theta,\varphi}$ for the scalar functional of its ordinary type, and put $Z_F=\int_X\mathrm{e}^F\,\mathrm{d}V$.

\begin{theorem}\label{thm:normalization-mass-quotient}
For every such pair and every bounded measurable $F$, the ordinary-type equation has a unique scalar and a solution unique modulo constants. Its functional and supporting set are
\[
 B_{\theta,\varphi}(F)=\log m_\varphi-\log Z_F,
 \qquad \mathcal A_{\theta,\varphi}(F)=\{\mathrm{e}^F\mathrm{d}V/Z_F\}.
\]
For two such pairs on the same $X,\mathrm{d}V$, possibly in different classes,
\begin{equation}\label{eq:normalization-mass-distance}
 \sup_{F\in L^\infty}|B_{\theta,\varphi}(F)-B_{\theta',\psi}(F)|
 =|\log m_\varphi-\log m_\psi|.
\end{equation}
Thus an invariant of these pairs is determined by the entire functional if and only if it is constant on every fiber of the mass map.
\end{theorem}
\begin{proof}
Existence follows from \cref{cor:big-model-lp-type}. Indeed, a K\"ahler manifold has bounded mass by \cref{lem:kahler-modification-bounded-mass}, and closed ordinary-type mass invariance gives the positive lower reference volume $m_\varphi$ in \eqref{eq:li-model-condition}. The classical closed prescribed-type theorem, recorded in \cref{sec:history-normalization}, gives uniqueness modulo constants under the stated model and small-unbounded-locus hypotheses.

Every ordinary-type competitor has mass $m_\varphi$. Integration of its equation forces $cZ_F=m_\varphi$, proving the scalar formula. As in \eqref{eq:normalization-closed-formula}, Jensen's inequality gives the stated support, and differentiation in every bounded direction makes it unique. Subtracting the scalar formulas gives \eqref{eq:normalization-mass-distance}. Finally $m_\varphi=\mathrm{d}V(X)\exp B_{\theta,\varphi}(0)$, so the functional and the mass determine each other. An invariant factors through either one exactly when it is constant on these fibers.
\end{proof}

For an ordinary divisorial type with a positive closed residual $\beta$, \cref{thm:ordinary-divisor-normalization} gives the same formula with mass $\int_X\beta^n$: bounded closed-product Stokes fixes this mass. This computation does not require a model representative for a real weighted sum of divisors. If $\beta_j\to\beta>0$ uniformly, continuity of the wedge polynomial gives convergence of the masses; the relative form bounds give the sharper inequalities $(1-\varepsilon_j)^n\int\beta^n\leq\int\beta_j^n\leq(1+\varepsilon_j)^n\int\beta^n$.

\begin{example}\label{ex:normalization-distance-projective}
Let $\theta$ be K\"ahler and let $\varphi,\psi$ be model references satisfying \eqref{eq:als-model-condition}, with positive non-pluripolar masses $m_\varphi,m_\psi$. The closed formula \eqref{eq:normalization-closed-formula} gives
\[
 \Delta(\varphi,\psi)=|\log m_\varphi-\log m_\psi|.
\]
For an explicit calculation, put $X=\mathbb P^n$ for an integer $n\geq1$ and let $\theta$ be its Fubini--Study form, normalized by $\int_X\theta^n=1$, and let $D$ be a hyperplane. Let $s_D$ be a defining section of $\mathcal O(1)$ with its Fubini--Study metric, and choose a rational $t\in(0,1)$. Then $a_D=\log|s_D|^2\leq0$ after rescaling, with
\[
 \ddc a_D=[D]-\theta,\qquad
 \theta+\ddc(ta_D)=(1-t)\theta+t[D].
\]
Put $\varphi_{t,D}=P_\theta[ta_D]$. As in \cref{ex:analytic-reference-normalization}, this is an admitted model reference of ordinary type $ta_D$. Off $D$ its original representative has Monge--Amp\`ere measure $(1-t)^n\theta^n$, and its non-pluripolar measure puts no mass on $D$.

\Needspace{5\baselineskip}
Invariance of mass under bounded changes in the closed background gives
\[
 m_{\varphi_{t,D}}=(1-t)^n,\qquad
 B_{\varphi_{t,D}}(G)-B_0(G)=n\log(1-t)\quad(G\in C(X,\mathbb R)).
\]
Thus the constant in \cref{prop:reference-normalization-comparison} is sharp.

\Needspace{10\baselineskip}
On $\mathbb P^1$, choose two distinct points $p,q$ for $D$. The corresponding references have distance zero, although one is unbounded at $p$ and bounded at $q$, while the other has the reverse behavior. They are not ordinarily equivalent. For every positive integer $m$,
\[
 \mathcal I(m\varphi_{t,p})=\mathcal O(-\lfloor mt\rfloor p),\qquad
 \mathcal I(m\varphi_{t,q})=\mathcal O(-\lfloor mt\rfloor q).
\]
These identities follow from \cref{lem:smooth-center-multiplier-ideal}. For $mt\geq1$ the ideals differ. Thus $\Delta$ does not separate ideal types either, and cannot replace the singularity topology sought in this chapter.
\end{example}

